% !TEX root = ../main.tex
% =====================================================================
%  CHƯƠNG 4 — THIẾT KẾ MẠCH BÙ NHIỆT
% ---------------------------------------------------------------------
%  Người phụ trách : [TODO: Họ tên - MSSV]
%  Tiền tố label   : c4      (fig:c4:..., tab:c4:..., eq:c4:..., sec:c4:...)
%  Thư mục ảnh     : image/chuong4/
%  File TLTK       : references/chuong4.bib   (key: c4_...)
%  Độ dài gợi ý    : 6–10 trang
%  Đáp ứng yêu cầu : (c).
%  Quy trình       : ảnh hưởng nhiệt (định lượng từ datasheet) -> phương pháp
%                    bù -> chọn cảm biến -> tính toán -> MÔ PHỎNG RIÊNG LẺ ở
%                    nhiều nhiệt độ (trước / sau bù) -> \nhanxet{}.
%  Giao tiếp       : Vị trí chèn mạch bù và dải tín hiệu phải khớp
%                    Bảng tab:c2:giao-tiep-khoi.
% =====================================================================
\chapter{Thiết kế mạch bù nhiệt}
\label{chap:c4}

% ---------------------------------------------------------------------
\section{Yêu cầu thiết kế}
\label{sec:c4:yeu-cau}

\TODO{Dải nhiệt độ làm việc, vị trí mạch bù trong sơ đồ khối (Hình~\ref{fig:c2:so-do-khoi}), sai số mục tiêu.}

% ---------------------------------------------------------------------
\section{Ảnh hưởng của nhiệt độ lên Load cell}
\label{sec:c4:anh-huong-nhiet}
% Định lượng trôi điểm không và trôi độ nhạy theo thông số datasheet,
% tính sai lệch điện áp ngõ ra tại các nhiệt độ biên.

\TODO{Tính độ trôi điểm không và độ nhạy theo nhiệt độ.}

\begin{equation}
    \Delta V_{\mathrm{zero}}(T) = \TODO{}
    \label{eq:c4:troi-diem-khong}
\end{equation}

% ---------------------------------------------------------------------
\section{Phương pháp bù nhiệt}
\label{sec:c4:phuong-phap}

\TODO{Trình bày phương án bù được chọn và lý do (so sánh với phương án khác nếu có).}

% ---------------------------------------------------------------------
\section{Lựa chọn cảm biến nhiệt độ}
\label{sec:c4:cam-bien}

\TODO{Lựa chọn cảm biến (RTD / NTC / IC), phương trình đặc tính, bảng thông số (có \textbackslash cite).}

% ---------------------------------------------------------------------
\section{Tính toán thiết kế mạch bù nhiệt}
\label{sec:c4:tinh-toan}

\TODO{Sơ đồ mạch, tính toán giá trị linh kiện để điện áp bù triệt tiêu độ trôi theo \eqref{eq:c4:troi-diem-khong}.}

% ---------------------------------------------------------------------
\section{Mô phỏng mạch bù nhiệt}
\label{sec:c4:mo-phong}

\TODO{Hình mô phỏng tại các nhiệt độ khác nhau, ví dụ \qty{0}{\celsius}, \qty{25}{\celsius}, \qty{50}{\celsius}.}

\begin{table}[H]
    \centering
    \caption{Điện áp ngõ ra trước và sau khi bù nhiệt}
    \label{tab:c4:ket-qua-bu-nhiet}
    \renewcommand{\arraystretch}{1.3}
    \begin{tabular}{|c|c|c|c|}
        \hline
        \multirow{2}{*}{\textbf{Nhiệt độ}} & \multirow{2}{*}{\textbf{Khối lượng}} &
        \multicolumn{2}{c|}{\textbf{Điện áp ngõ ra (\unit{\volt})}} \\ \cline{3-4}
        & & \textbf{Trước bù} & \textbf{Sau bù} \\ \hline
        \qty{0}{\celsius}  & \qty{0}{\kilo\gram} & \TODO{} & \TODO{} \\ \hline
        \qty{25}{\celsius} & \qty{0}{\kilo\gram} & \TODO{} & \TODO{} \\ \hline
        \qty{50}{\celsius} & \qty{0}{\kilo\gram} & \TODO{} & \TODO{} \\ \hline
        \qty{0}{\celsius}  & \qty{1}{\kilo\gram} & \TODO{} & \TODO{} \\ \hline
        \qty{25}{\celsius} & \qty{1}{\kilo\gram} & \TODO{} & \TODO{} \\ \hline
        \qty{50}{\celsius} & \qty{1}{\kilo\gram} & \TODO{} & \TODO{} \\ \hline
    \end{tabular}
\end{table}

\nhanxet{\TODO{Mức độ giảm sai lệch nhờ mạch bù, so sánh với giá trị tại \qty{25}{\celsius}.}}

% ---------------------------------------------------------------------
\section{Tổng kết}
\label{sec:c4:tong-ket}

\TODO{Tóm tắt kết quả mạch bù nhiệt, dải nhiệt độ đạt yêu cầu.}
